\section*{Kurzfassung}
\label{sec:kurzfassung}
\addcontentsline{toc}{section}{\nameref{sec:kurzfassung}}

Die Syntax von Markdown ist klein und gut lesbar, weshalb sie sich zum Verfassen von
Dokumentation eignet: Da es sich um reinen Text handelt, kann er in jedem beliebigen Texteditor geöffnet und
bearbeitet werden, und Änderungen daran lassen sich mühelos von Versionsverwaltungssystemen
nachverfolgen.
Für das Schreiben von Lehrmaterial ist Markdown allerdings nur bedingt geeignet -- nicht, weil es
selbst zu kurz griffe, sondern weil das Lernmanagementsystem Moodle es nicht unterstützt.
In Moodle können Inhalte als H5P-Pakete in einem browserbasierten Editor erstellt werden.
Das bedeutet jedoch, dass in Markdown verfasstes Material von Hand portiert wird, und diese
Portierung wiederholt sich bei jeder Änderung der Quelle.
H5P-Pakete lassen sich für verschiedene Inhaltstypen erzeugen, von denen Interactive Book
Markdown-Dateien am ehesten gerecht wird: Dieser Typ gliedert Inhalte in Kapitel, wobei jedes
Kapitel aus einzelnen Inhaltsblöcken besteht.
Diese Blöcke stammen aus einer kleinen Auswahl von Inhaltstypen, die festlegen, was
worin verschachtelt werden darf.
Manche Markdown-Konstrukte -- etwa ein
Bild in einer Tabellenzelle -- haben daher keine unmittelbare Entsprechung.
Die vorliegende Arbeit lotet die Grenzen des H5P-Inhaltstyps Interactive Book aus und stellt
md2h5p vor, einen Konverter, der einen Ordner mit Markdown-Dateien in ein H5P Interactive
Book überführt.
Er übersetzt jede Markdown-Datei in ein Kapitel, findet für Markdown-Konstrukte, die H5P nicht
direkt unterstützt, geeignete Alternativen und gibt Warnungen für Markdown aus, das sich nur in
verminderter Form -- oder gar nicht -- darstellen lässt.
Eine Erweiterung ergänzt Single- und Multiple-Choice-Fragen, die als eingezäunte Codeblöcke in
Markdown geschrieben werden, und ein Merge-Schritt verbindet eine neue Konvertierung mit einem
Paket, das zwischenzeitlich im H5P-Editor bearbeitet wurde.
Lehrende können somit weiterhin in Markdown schreiben, ihr Material auch nach der Veröffentlichung
überarbeiten und ihren Studierenden dennoch Inhalte bereitstellen, die die Plattform einbetten kann.
